% $Id:abstract.tex  $
% !TEX root = main.tex

\chapter{Abstract}

Connected vehicles send safety messages directly to each other over vehicle-to-vehicle (V2V) links, which can
reuse the uplink channels of the users of the cellular network at the cost of interfering with them. Choosing
the channel and the transmit power of each pair of vehicles is a sequential decision problem that has been
addressed with multi-agent reinforcement learning. In the scheme of Li et al.~\cite{li+22fedmarl}, the base
station of the cell periodically averages the neural networks of all the agents (FedAvg), and all the agents
learn from a single shared reward that the base station computes.

This thesis revisits these two choices. First, it proposes a decentralized aggregation in which the base
station takes no part in training: each vehicle averages its network only with those of the vehicles within
\SI{300}{\meter} (gossip aggregation), and the groups of neighbors change with the positions of the
vehicles. Second, it proposes an individual reward, which judges each agent by what its own action can
change instead of by the sum of what all the agents do.

We implemented the complete system and evaluated it in simulation with 5 to 30 pairs of vehicles, judging
every method by the shared reward. The individual reward improved the result in almost every paired
comparison, and with it, gossip beat FedAvg in the middle range of the number of pairs; without it, gossip
brought no consistent gain. Together, the two changes beat our implementation of the original scheme in all
paired comparisons. These gains have a limit: a simple fixed rule that uses only local information still
beats every learning method.

\medskip
\noindent\textbf{Keywords:} vehicle-to-vehicle communications; radio resource allocation; multi-agent
reinforcement learning; decentralized federated learning; gossip; reward design.

\endinput
